\documentclass[11pt]{article}

\usepackage[letterpaper,margin=0.82in,headheight=15pt]{geometry}
\usepackage[T1]{fontenc}
\usepackage{lmodern}
\usepackage{microtype}
\usepackage{amsmath,amssymb,bm,mathtools}
\usepackage{booktabs,longtable,tabularx,array}
\usepackage{enumitem}
\usepackage{xcolor}
\usepackage{fancyhdr}
\usepackage{hyperref}
\usepackage[nameinlink,noabbrev]{cleveref}

\definecolor{navy}{HTML}{17324D}
\definecolor{teal}{HTML}{147D92}
\definecolor{orange}{HTML}{D47A27}
\definecolor{softblue}{HTML}{EEF4F7}
\definecolor{softgray}{HTML}{F3F5F6}
\definecolor{darkgray}{HTML}{46545E}

\hypersetup{
  colorlinks=true,
  linkcolor=navy,
  citecolor=teal,
  urlcolor=teal,
  pdftitle={An algebraic and geometric architecture for a predictive spin-1 light-front model},
  pdfauthor={DeuteronWigner project}
}

\pagestyle{fancy}
\fancyhf{}
\fancyhead[L]{\small\color{darkgray}Predictive spin-1 light-front architecture}
\fancyhead[R]{\small\color{darkgray}Algebraic and geometric formulation}
\fancyfoot[L]{\small\color{darkgray}27 July 2026}
\fancyfoot[R]{\small\color{darkgray}\thepage}
\renewcommand{\headrulewidth}{0.4pt}

\setlist[itemize]{leftmargin=1.5em,itemsep=2pt,topsep=3pt}
\setlist[enumerate]{leftmargin=1.7em,itemsep=2pt,topsep=3pt}
\setlength{\parindent}{0pt}
\setlength{\parskip}{5pt}
\setlength{\emergencystretch}{2em}
\renewcommand{\arraystretch}{1.15}

\newcommand{\kT}{\bm{k}_{T}}
\newcommand{\bT}{\bm{b}_{T}}
\newcommand{\DT}{\bm{\Delta}_{T}}
\newcommand{\dd}{\mathrm{d}}
\newcommand{\Tr}{\operatorname{Tr}}
\newcommand{\Hil}{\mathcal H}
\newcommand{\Alg}{\mathcal A}
\newcommand{\Rep}{\mathcal V}
\newcommand{\Path}{\mathrm{Path}}
\newcommand{\Obj}{\operatorname{Obj}}
\newcommand{\Hom}{\operatorname{Hom}}
\newcommand{\End}{\operatorname{End}}
\newcommand{\LF}{\mathrm{LF}}
\newcommand{\TMD}{\mathrm{TMD}}
\newcommand{\GTMD}{\mathrm{GTMD}}
\newcommand{\CP}{\mathrm{CP}}

\newcolumntype{Y}{>{\raggedright\arraybackslash}X}
\newcolumntype{L}[1]{>{\raggedright\arraybackslash}p{#1}}

\title{\vspace{-1.0cm}
{\color{navy}\bfseries An Algebraic and Geometric Architecture\\[2mm]
for a Genuinely Predictive Spin-1 Light-Front Model}\\[4mm]
\large Physical motivation, mathematical formulation, computational role,
and limits of the proposal}
\author{DeuteronWigner project}
\date{27 July 2026}

\begin{document}
\maketitle

\begin{center}
\begin{minipage}{0.93\textwidth}
\colorbox{softblue}{\parbox{0.96\textwidth}{
\textbf{Purpose.}
This note explains a proposed mathematical architecture for the next-level
model identified in Section 15 of
\texttt{references/model\_construction\_note.tex}.  The proposal is not that
topology supplies missing QCD dynamics.  It is that graded Hilbert spaces,
representation-theoretic intertwiners, gauge-bundle parallel transport,
symplectic phase-space geometry, and convex positive-state geometry can make
the physically required relations structural rather than optional.  A more
limited chain-complex construction is proposed as an auditable bookkeeping
device for Fock-sector transitions and double counting.
}}
\end{minipage}
\end{center}

\begin{abstract}
The present DeuteronWigner model is a constrained phenomenological synthesis:
it preserves a common correlator architecture and many exact constraints, but
important nucleon, gauge-link, tensor, gluon, and non-nucleonic components are
still informed by separately fitted or modeled inputs.  A genuinely
predictive replacement must instead generate many observables from a small
set of shared microscopic interactions and states.  We propose organizing
that calculation as a gauge-covariant graded light-front tensor category.
The dynamical core is a regulated light-front Hamiltonian on flavor-, color-,
spin-, orbital-, and Fock-graded Hilbert spaces.  Symmetry intertwiners define
allowed couplings.  Gauge links are parallel transports in an \(SU(3)\)
bundle and form a representation of a path groupoid.  GTMDs are common
operator matrix elements whose TMD, GPD, form-factor, current, and OAM
reductions must form commuting diagrams.  Transverse phase space is treated
symplectically, parton--target correlators inhabit a convex positive cone, and
nucleon-to-deuteron composition is represented, where applicable, by
completely positive amplitude maps.  A filtered family of regulated theories
defines convergence and renormalization tests.  We describe why each layer is
physically justified, what it prevents, how it could be implemented, and
which stronger topological claims should not be made.
\end{abstract}

\tableofcontents

\section{The problem this formulation is meant to solve}

The desired advance is a change in the source of prediction.  In the present
model, many ingredients are physically informed and assembled under a common
correlator contract,
\begin{equation}
\{\hbox{PDF/TMD fits, wave functions, response models, phase models}\}
\longrightarrow W_D
\longrightarrow \{F_i^D\}.
\label{eq:assembled}
\end{equation}
This is scientifically useful, but correlations among different \(F_i^D\)
are only as fundamental as the assumptions used in the assembly.  The
next-level direction should instead be
\begin{equation}
\{H_{\LF},\mathcal R,\theta\}
\longrightarrow
\{|\Psi_N\rangle,|\Psi_D\rangle,U_\gamma\}
\longrightarrow W_D
\longrightarrow \{F_i^D,\mathcal O_{\rm process}\},
\label{eq:generative}
\end{equation}
where \(H_{\LF}\) is a regulated and renormalized microscopic Hamiltonian,
\(\mathcal R\) denotes the regulator and matching prescription, and
\(\theta\) is a comparatively small set of shared interaction parameters.

The mathematical problem is compositional.  A state carries simultaneously
Fock sector, flavor, color, parton helicity, target helicity, OAM, gauge-link,
and nuclear-component information.  Every contraction must respect exact
selection rules.  Every projected TMD must remain traceable to the same
amplitudes.  Every truncation must have a convergence meaning.  A flat array
of independently parameterized functions is a poor representation of this
structure.

\subsection{Design criterion}

We will regard a mathematical structure as useful only if it does at least
one of the following:
\begin{enumerate}
\item makes an exact symmetry or admissibility condition automatic;
\item prevents an inconsistent or double-counted composition;
\item exposes a controlled approximation and its convergence direction;
\item creates testable relations among observables derived from a common
state;
\item reduces redundant parameterization without suppressing a physical
degree of freedom.
\end{enumerate}
Abstract terminology that meets none of these criteria should not enter the
production model.

\section{Graded light-front Hilbert space}

\subsection{Sector and internal grading}

The state space should be a direct sum
\begin{equation}
\Hil_{\LF}
=
\bigoplus_{\nu\in\mathcal F}\Hil_\nu,
\qquad
\nu=(n_q,n_{\bar q},n_g;B,Q,S,\ldots),
\label{eq:graded-hilbert}
\end{equation}
over a set \(\mathcal F\) of retained Fock sectors.  Each sector has the
form
\begin{equation}
\Hil_\nu
=
L^2(\mathcal M_\nu,\dd\mu_\nu)
\otimes\Rep_{\rm flavor}
\otimes\Rep_{\rm color}
\otimes\Rep_{\rm spin}
\otimes\Rep_{\rm OAM},
\label{eq:sector-space}
\end{equation}
where
\begin{equation}
\mathcal M_\nu=
\left\{(x_i,\bm k_{iT})\,\middle|\,
x_i>0,\quad
\sum_i x_i=1,\quad
\sum_i\bm k_{iT}=0
\right\}.
\label{eq:momentum-manifold}
\end{equation}
The measure contains the conventional light-front phase-space factors and
the constraint delta functions.

This formulation gives support and momentum conservation a geometric
location: amplitudes are defined on the momentum simplex and transverse
constraint surface.  Flavor, color, spin, and OAM are not metadata but
representation factors acted on by explicit operators.

\subsection{Block Hamiltonian and Fock transitions}

The mass eigenproblem is
\begin{equation}
H_{\LF}|\Psi_h(P,\Lambda)\rangle
=M_h^2|\Psi_h(P,\Lambda)\rangle,
\qquad
H_{\LF}=P^+P^- -\bm P_T^2.
\label{eq:lf-eigen}
\end{equation}
With sector projectors \(P_\nu\), define
\begin{equation}
H_{\nu\mu}=P_\nu H_{\LF}P_\mu.
\end{equation}
The diagonal blocks contain kinetic and sector-preserving interactions;
off-diagonal blocks generate emissions, absorption, and pair creation:
\begin{equation}
H_{\LF}\sim
\begin{pmatrix}
H_{qqq} & V_{qqq,qqqg} & V_{qqq,qqqq\bar q}&\cdots\\
V^\dagger_{qqq,qqqg}&H_{qqqg}&V_{qqqg,qqqq\bar q}&\cdots\\
V^\dagger_{qqq,qqqq\bar q}&V^\dagger_{qqqg,qqqq\bar q}
&H_{qqqq\bar q}&\cdots\\
\vdots&\vdots&\vdots&\ddots
\end{pmatrix}.
\label{eq:block-hamiltonian}
\end{equation}
This is the natural setting for basis light-front quantization (BLFQ),
Hamiltonian effective field theory, or a compatible continuum bound-state
solution mapped to light-front amplitudes \cite{Vary2009}.

\subsection{Why the grading matters}

The grading provides explicit answers to questions that are otherwise
hidden:
\begin{itemize}
\item Which sea distribution is generated by which pair-creation sector?
\item Which gluon TMD receives support from \(qqqg\) rather than an assigned
gluon profile?
\item Which spin--orbit observable requires interference between which
\(L_z\) sectors?
\item Which counterterm compensates the integration-out of a sector?
\item Does an apparent prediction stabilize as sectors are enlarged?
\end{itemize}
It also prevents a nucleon gluon input and an explicit \(qqqg\) sector from
being included without an explicit matching or subtraction rule.

\section{Representation theory as the coupling grammar}

\subsection{Symmetry representations}

At a regulated stage, the relevant exact or controlled symmetry action is
organized by
\begin{equation}
G=
SU(3)_{\rm color}
\times SU(2)_{\rm isospin}
\times SU(2)_{\rm spin}
\times SO(2)_{L_z},
\label{eq:group-product}
\end{equation}
supplemented by discrete parity, time-reversal, and charge-conjugation maps.
The \(SU(2)_{\rm spin}\) factor must be understood with the appropriate
light-front spin and Melosh/Wigner rotations; full Poincare covariance is a
continuum-limit requirement rather than something guaranteed by writing
\(SU(2)\).

An interaction tensor is an intertwiner
\begin{equation}
C:\Rep_{R_1}\otimes\cdots\otimes\Rep_{R_n}\longrightarrow\Rep_{R_{\rm out}},
\qquad
C\,D_{\rm in}(g)=D_{\rm out}(g)\,C.
\label{eq:intertwiner}
\end{equation}
Only invariant or correctly covariant tensors are parameterized.  Selection
rules are therefore restrictions on the space of maps, not penalties added
after fitting.

\subsection{Spin-1 irreducible density structure}

For a spin-1 target,
\begin{equation}
1\otimes1^*=0\oplus1\oplus2.
\label{eq:spin1-decomp}
\end{equation}
Accordingly, the density operator decomposes into scalar, vector, and
quadrupole pieces,
\begin{equation}
\rho_D
=\rho_D^{(0)}
+\sum_{m=-1}^{1}S_m\,T_m^{(1)}
+\sum_{m=-2}^{2}Q_m\,T_m^{(2)}.
\label{eq:density-irrep}
\end{equation}
The familiar \(U,L,T,LL,LT,TT\) labels are basis-dependent real
combinations of these irreducible components.  Working first in the
spherical-tensor basis has concrete benefits:
\begin{itemize}
\item rotations are implemented by Wigner \(D\)-matrices;
\item vector and tensor responses cannot be silently conflated;
\item pure-\(S\), \(S\)--\(D\), and additional OAM limits have explicit
Clebsch--Gordan selection rules;
\item the Cartesian TMD basis is produced by a fixed invertible map.
\end{itemize}

\subsection{Flavor and controlled isospin}

Proton and neutron states should inhabit an isospin multiplet, but physical
operators may include controlled breaking:
\begin{equation}
H_{\LF}=H_0+\delta H_{m_d-m_u}+\delta H_{\rm QED}
+\delta H_{\rm nuclear}.
\label{eq:isospin-breaking}
\end{equation}
Exact isospin is recovered when the last three terms are disabled.  It is
not implemented by assigning identical \(u\) and \(d\) functions.  The
representation structure therefore permits both exact interchange tests and
realistic charge-symmetry breaking.

\subsection{Color invariants and gluon T-odd sectors}

The two three-adjoint invariant tensors
\begin{equation}
f^{abc},\qquad d^{abc}
\end{equation}
define inequivalent color contractions.  They should be different
intertwiner channels in the operator graph.  This prevents a common scalar
``gluon T-odd amplitude'' from being copied into both sectors and allows a
hard process to select its own linear combination.

\section{A symmetry-preserving tensor-network realization}

\subsection{What the network represents}

A tensor network is proposed as a sparse computational representation of
the amplitude and operator contractions, not as an assumption that QCD is
one-dimensional or has low entanglement everywhere.  Edges carry
\begin{equation}
e=(\nu,R_c,R_I,j,\lambda,L_z,\alpha),
\label{eq:edge-label}
\end{equation}
and nodes are physical maps: Hamiltonian vertices, spin couplers,
wave-function coefficients, current insertions, Wilson-line vertices, or
nuclear Fock-sector couplers.

For example, a schematic deuteron amplitude is
\begin{equation}
|\Psi_D^\Lambda\rangle
=
\mathcal C_D^\Lambda\!
\left[
|\Psi_p\rangle\otimes|\Psi_n\rangle
\oplus|\Psi_{NN\pi}\rangle
\oplus|\Psi_{\Delta\Delta}\rangle
\oplus|\Psi_{6q}\rangle+\cdots
\right].
\label{eq:deuteron-network}
\end{equation}
The \(NN\) coupling contains maps
\begin{equation}
C_{D,L}:
\left(\tfrac12\otimes\tfrac12\right)_{S=1}
\otimes L
\longrightarrow J=1,
\qquad L=0,2,\ldots .
\label{eq:deuteron-intertwiner}
\end{equation}

\subsection{What it buys us}

The network makes shared microscopic origin inspectable.  A parameter lives
on one Hamiltonian vertex or state tensor and propagates to every observable
that uses it.  It cannot become an independent coefficient for each TMD.
Automatic differentiation through contractions can produce correlated
sensitivities
\begin{equation}
J_{ia}=\frac{\partial\mathcal O_i}{\partial\theta_a},
\qquad
\operatorname{Cov}(\mathcal O)
\simeq J\,\operatorname{Cov}(\theta)\,J^T,
\label{eq:network-covariance}
\end{equation}
while retaining the exact symmetry blocks.

\subsection{What it does not buy us}

A tensor network does not make the interaction correct.  Aggressive
low-rank truncation may discard precisely the OAM, color, or long-range
correlations of interest.  Bond-dimension convergence must therefore be
treated on the same footing as basis and Fock-sector convergence.  The
network is a representation and contraction strategy, not a phenomenological
replacement for \(H_{\LF}\).

\section{Gauge links as bundle parallel transport}

\subsection{Connection and holonomy}

Let \(P(M,SU(3))\) be the color principal bundle over the relevant spacetime
domain and
\begin{equation}
A=A_\mu^a t^a\,\dd x^\mu
\end{equation}
its connection.  The gauge link along an oriented path \(\gamma:a\to b\) is
\begin{equation}
U_\gamma[b,a]
=\mathcal P\exp\left[-ig\int_\gamma A\right].
\label{eq:holonomy}
\end{equation}
Parallel transport obeys the exact algebra
\begin{equation}
U_{\gamma_2\circ\gamma_1}
=U_{\gamma_2}U_{\gamma_1},
\qquad
U_{\gamma^{-1}}=U_\gamma^\dagger .
\label{eq:path-algebra}
\end{equation}
Quarks use the fundamental representation and gluons the adjoint
representation.  Staple direction, transverse closure, rapidity regulator,
and representation are properties of the operator itself.

\subsection{The path groupoid}

The admissible gauge-link paths form a groupoid:
\begin{align}
\Obj(\Path)&=\{\hbox{operator endpoints}\},\\
\Hom_{\Path}(a,b)&=\{\gamma:a\to b\},\\
\gamma^{-1}&:b\to a.
\end{align}
The Wilson-line assignment is a representation
\begin{equation}
\mathcal U_R:\Path\longrightarrow\operatorname{End}(\Rep_R),
\qquad
\gamma\longmapsto U_\gamma^{(R)}.
\label{eq:path-representation}
\end{equation}
Future- and past-pointing staples are distinct morphisms.  Time reversal maps
their classes into one another, so the familiar T-even/T-odd relations are
implemented as transformations of operators:
\begin{align}
F_{\rm even}^{[\gamma_+]}&=F_{\rm even}^{[\gamma_-]},\\
F_{\rm odd}^{[\gamma_+]}&=-F_{\rm odd}^{[\gamma_-]},
\label{eq:link-reversal}
\end{align}
subject to the relevant conjugations and conventions.

\subsection{Where the dynamics enters}

The curvature
\begin{equation}
F=\dd A-igA\wedge A
\label{eq:curvature}
\end{equation}
controls local path dependence.  Expanding the holonomy,
\begin{equation}
U_\gamma
=1-ig\int_\gamma A
-g^2\int_{\gamma}A\,A+\cdots ,
\label{eq:holonomy-expansion}
\end{equation}
shows that absorptive phases and naive-T-odd functions must arise from
intermediate states, pole prescriptions, and rescattering dynamics.  The
path groupoid organizes which phase is allowed; it does not determine the
phase numerically.  That requires the Hamiltonian gauge field and the
factorization-appropriate eikonal limit.

\subsection{Why this is geometric but not topological QCD}

For a flat connection, holonomy may depend only on homotopy class.  QCD
generically has \(F_{\mu\nu}\ne0\), so the Wilson line depends on the detailed
path geometry.  It is therefore incorrect to call a Sivers or
Boer--Mulders function a topological invariant.  Topology can classify link
connectivity and closed cycles; the non-Abelian connection supplies their
physical values.  This distinction is essential if the construction is to
remain falsifiable.

\section{GTMDs as common morphisms and commuting reductions}

\subsection{Operator matrix element}

For an operator label \(\Gamma\) and link path \(\gamma\), define
\begin{equation}
\mathcal W_{\Gamma,\gamma}:\Hil_h\longrightarrow\Hil_h.
\end{equation}
Its off-forward matrix element is
\begin{equation}
W_{\Lambda'\Lambda}^{[\Gamma,\gamma]}
(x,\kT,\DT)
=
\langle P',\Lambda'|
\mathcal W_{\Gamma,\gamma}(x,\kT,\DT)
|P,\Lambda\rangle .
\label{eq:gtmd-morphism}
\end{equation}
A wave-function calculation evaluates \cref{eq:gtmd-morphism} as overlaps of
the same Fock amplitudes, including diagonal and, where required,
off-diagonal sector terms.  GTMDs supply a natural common ancestor for TMD
and GPD limits \cite{Meissner2008}.

\subsection{Reduction maps}

Define schematic linear maps
\begin{align}
\mathrm T[W]&=W|_{\DT=0},&
\mathrm G[W]&=\int\dd^2\kT\,W,\nonumber\\
\mathrm F[W]&=\int\dd x\,\dd^2\kT\,W,&
\mathrm Q_\Gamma[W]&=\hbox{local-current reduction}.
\label{eq:reduction-maps}
\end{align}
The predictive architecture requires compatible reductions to commute:
\begin{equation}
\mathrm F\circ\mathrm T[W]
=
\mathrm Q_\Gamma[W]
\end{equation}
where the regulator, matching scheme, and operator definitions make the
identity applicable.  Similarly, nuclear composition \(\mathrm N\) and
projection \(\mathrm P_i\) should satisfy
\begin{equation}
\mathrm P_i\circ\mathrm N[W_N]
\stackrel{?}{=}
\mathrm N_i\circ\mathrm P[W_N].
\label{eq:nuclear-commute}
\end{equation}
Failure of \cref{eq:nuclear-commute} is not automatically an error: it may
show that projection-dependent many-body currents or coherent terms are
missing.  The commuting diagram is therefore both an architecture and a
diagnostic.

\subsection{Why this is more predictive}

An independently fitted \(f_1\), \(g_1\), \(h_1\), Sivers function, and
pretzelosity need not satisfy any common off-forward, current, or OAM
relation.  If all are reductions of \cref{eq:gtmd-morphism}, their
correlations follow from the state and operator.  A failure against a
withheld observable then challenges the shared amplitudes rather than one
isolated curve.

\section{Symplectic geometry of transverse phase space}

\subsection{Phase space and moment map}

The transverse Wigner variables inhabit
\begin{equation}
T^*\mathbb R^2
=\{(\bT,\kT)\},
\end{equation}
with canonical symplectic form
\begin{equation}
\omega
=\dd b_x\wedge\dd k_x+\dd b_y\wedge\dd k_y.
\label{eq:symplectic-form}
\end{equation}
The \(SO(2)\) rotation action has moment map
\begin{equation}
\mu_{L_z}(\bT,\kT)
=b_xk_y-b_yk_x.
\label{eq:moment-map}
\end{equation}
Thus an OAM expectation is a phase-space moment
\begin{equation}
\langle L_z\rangle
=
\int\dd x\,\dd^2\bT\,\dd^2\kT\,
\mu_{L_z}(\bT,\kT)\,
\rho_W(x,\bT,\kT;\gamma),
\label{eq:oam-wigner}
\end{equation}
with an operator definition that retains its Wilson-line dependence
\cite{Lorce2012}.

\subsection{Transverse rank as representation weight}

Functions \(e^{im\phi}\) carry irreducible \(SO(2)\) weight \(m\).
Transverse-rank tensors are therefore organized by angular harmonics rather
than by visually assigned powers of \(k_T\).  A matrix element obeys
\begin{equation}
\Delta L_z+\Delta\lambda_{\rm parton}
+\Delta\Lambda_{\rm target}=m_{\rm operator}.
\label{eq:oam-selection}
\end{equation}
Equation~\eqref{eq:oam-selection} identifies which partial-wave
interferences can contribute to worm gears, pretzelosity, tensor TMDs, and
linearly polarized gluon distributions.  It also supplies controlled
zero tests: a required interference vanishes when either participating
amplitude is removed.

\subsection{Caution about Wigner functions}

The Wigner distribution is a quasiprobability and need not be pointwise
positive.  Symplectic geometry organizes its transformations and moments;
it must not be confused with the positive density-matrix cone discussed
next.

\section{Convex geometry and positivity by construction}

\subsection{The correlator cone}

At fixed kinematics, the combined parton--target helicity correlator must
be Hermitian and positive semidefinite in the applicable forward
probability interpretation:
\begin{equation}
\Phi(x,\kT)\in\CP_n
\equiv\{X=X^\dagger\mid X\succeq0\}.
\label{eq:positive-cone}
\end{equation}
Named TMDs are linear coordinates on the image of this cone under a
projection map.  Their allowed region is therefore a convex spectrahedral
set, not a Cartesian box of independent error bars.

The most physical factorization is
\begin{equation}
\Phi_{\alpha\beta}
=\sum_X A_{\alpha X}A_{\beta X}^*,
\label{eq:amplitude-factor}
\end{equation}
where \(X\) labels intermediate or spectator states.  Numerically one may
also use a Cholesky-like form
\begin{equation}
\Phi=LL^\dagger.
\label{eq:cholesky}
\end{equation}
Hermiticity and positivity then hold for every parameter value.  Principal
minors yield Soffer-like and spin-1 tensor bounds.  This density-matrix
view is already the basis of known spin-1 gluon TMD positivity constraints
\cite{Cotogno2017}.

\subsection{Why clipping is unacceptable}

If independently generated TMDs are assembled and negative eigenvalues are
clipped, the model changes in a nonlinear and untraceable way.  In
\cref{eq:amplitude-factor}, a positivity failure instead means that the
claimed amplitudes or projection map are inconsistent.  Uncertainty samples
remain inside the physical cone without erasing correlations among
different functions.

\subsection{Information geometry}

For normalized positive states, distinguishability may be quantified by
the quantum Fisher metric or relative entropy.  This has a concrete
possible use in experiment design: identify observables that separate two
microscopic mechanisms rather than merely enlarging an already constrained
direction.  It should be introduced only after the physical density
operator and likelihood are fixed.

\section{Nucleon-to-deuteron composition as an amplitude map}

\subsection{Completely positive impulse composition}

For an incoherent retained subsystem, the nucleon-to-deuteron map can be
written
\begin{equation}
\mathcal E_D(\rho_N)
=
\sum_\alpha K_\alpha\rho_NK_\alpha^\dagger,
\label{eq:kraus}
\end{equation}
where \(\alpha\) includes spectator helicity, partial wave, recoil, and
retained nuclear channel.  If the map is trace preserving,
\begin{equation}
\sum_\alpha K_\alpha^\dagger K_\alpha=I.
\label{eq:trace-preserving}
\end{equation}
If probability flows to \(NN\pi\), \(\Delta\Delta\), or other explicit
sectors, the corresponding maps must be added before imposing total
normalization.

The value of \cref{eq:kraus} is that positivity is inherited from the
nucleon correlator.  Tensor polarization is generated by the angular and
spin structure of the \(K_\alpha\), rather than by multiplying an
unpolarized answer by an arbitrary tensor factor.

\subsection{Limits of the quantum-channel picture}

Coherent shadowing, phases between Fock sectors, and many-body currents
must be computed at the amplitude level before an environmental trace.
An apparently non-completely-positive reduced approximation may indicate
that a coherent degree of freedom was traced out inconsistently.  We should
not force every nuclear mechanism into an incoherent Kraus map.  Instead,
the full amplitude is primary and the channel representation is used where
its assumptions are satisfied.

\section{Chain complexes for composition and double counting}

\subsection{A deliberately limited topological layer}

Let \(C_n\) be the vector space generated by retained configurations with
\(n\) resolved transitions relative to a reference sector.  A boundary map
\begin{equation}
\partial_n:C_n\longrightarrow C_{n-1},
\qquad
\partial_{n-1}\partial_n=0,
\label{eq:chain-complex}
\end{equation}
records the endpoints and signs of sector-transition sequences.  Examples
include
\[
NN\leftrightarrow NN\pi,\qquad
qqq\leftrightarrow qqqg,\qquad
qqq\leftrightarrow qqqq\bar q.
\]

This layer can support:
\begin{itemize}
\item automated detection of the same pion dressing represented both inside
a nucleon input and as an explicit \(NN\pi\) sector;
\item association of effective operators and counterterms with the sectors
that were integrated out;
\item classification of closed transition cycles requiring interference or
unitarity checks;
\item provenance paths from a final observable back to its amplitude
components.
\end{itemize}

The homology
\begin{equation}
H_n=\ker\partial_n/\operatorname{im}\partial_{n+1}
\label{eq:homology}
\end{equation}
would classify independent cycles not reducible to already represented
higher transitions.  This is initially an audit construction, not a claim
that physical amplitudes depend only on a homology class.

\subsection{Why category language is also useful}

Objects can represent typed state/operator spaces; morphisms represent
Hamiltonian evolution, Wilson transport, projection, matching, evolution,
or nuclear composition.  Composition is allowed only when domains,
codomains, schemes, link classes, and quantum numbers match.  This provides
a mathematically precise version of the project's fail-closed typed
interfaces.

We do not need sophisticated category theory to implement this.  Sparse
block tensors plus checked domain/codomain metadata are sufficient.  The
categorical viewpoint is valuable because it states what must compose and
which diagrams must commute.

\section{Regulator, renormalization, and convergence as a filtration}

\subsection{Nested regulated spaces}

Define a family of calculations
\begin{equation}
\Hil_1\hookrightarrow\Hil_2\hookrightarrow\cdots,
\qquad
\Hil_n=\Hil(\Lambda_n,N_{\max,n},\mathcal F_n),
\label{eq:filtration}
\end{equation}
that enlarges ultraviolet resolution, basis size, and Fock content.  Each
level has
\begin{equation}
H^{(n)}=H_0^{(n)}+\sum_a c_a^{(n)}O_a^{(n)}.
\end{equation}
Matching maps
\begin{equation}
R_{n+1\to n}:\Alg_{n+1}\longrightarrow\Alg_n
\label{eq:matching-map}
\end{equation}
relate effective operators across truncations.  For calibration observables,
\begin{equation}
R_{n+1\to n}[\mathcal O^{(n+1)}]
=\mathcal O^{(n)}+\delta_n,
\qquad
\delta_n\to0
\label{eq:projective-consistency}
\end{equation}
within an estimated truncation law.

\subsection{Acceptance meaning}

A smooth result at one \(N_{\max}\) is not a prediction.  A predictive
claim requires stability under:
\begin{itemize}
\item basis enlargement;
\item regulator variation with counterterm refitting;
\item addition of the next physically required Fock sector;
\item increased tensor-network bond or factorization rank;
\item quadrature and grid refinement.
\end{itemize}
The correlated spread across these directions is a model-discrepancy
component, not interchangeable with posterior parameter uncertainty.

\section{Inference and falsifiability}

\subsection{Shared microscopic parameters}

Let \(\theta\) contain masses, couplings, regulator-dependent counterterms,
and a restricted number of effective interactions.  The posterior is
\begin{equation}
p(\theta,\delta_{\rm EFT}\mid D_{\rm cal})
\propto
p(D_{\rm cal}\mid\theta,\delta_{\rm EFT})
p(\theta)p(\delta_{\rm EFT}).
\label{eq:posterior}
\end{equation}
The calibration data may include masses, elastic currents, selected
nucleon PDFs/charges, deuteron static properties, and a limited set of
process observables.  It must not include a free normalization and
transverse width for every TMD.

The predictive distribution for withheld data is
\begin{equation}
p(D_{\rm hold}\mid D_{\rm cal})
=\int\dd\theta\,\dd\delta_{\rm EFT}\,
p(D_{\rm hold}\mid\theta,\delta_{\rm EFT})
p(\theta,\delta_{\rm EFT}\mid D_{\rm cal}).
\label{eq:posterior-predictive}
\end{equation}
A failed withheld family must revise the state, interaction, current, or
truncation model.  Adding a coefficient directly to the failed TMD would
return us to the patchwork construction.

\subsection{The prediction hierarchy}

The strongest test is cross-channel:
\begin{enumerate}
\item calibrate a common Hamiltonian using spectroscopy, currents, and a
restricted set of partonic data;
\item predict GTMD marginals not used in calibration;
\item predict link-reversal or color-channel relations;
\item predict tensor-polarized and gluon observables;
\item propagate without retuning into deuteron observables.
\end{enumerate}
The framework is worthwhile only if it makes such tests possible.

\section{Software realization}

\subsection{Core objects}

\begin{longtable}{L{0.28\textwidth}L{0.64\textwidth}}
\toprule
Object & Responsibility\\
\midrule
\endhead
\texttt{SectorSpace} &
Momentum manifold, basis, Fock grade, representations, regulator, and
normalization.\\
\texttt{Intertwiner} &
Sparse invariant tensor with explicit domain, codomain, multiplicity, and
symmetry tests.\\
\texttt{LFHamiltonian} &
Hermitian block operator, counterterms, sector-changing vertices, and
eigensolver adapter.\\
\texttt{FockState} &
Normalized amplitudes, phase convention, member identity, and truncation
metadata.\\
\texttt{PathGroupoid} and \texttt{WilsonTransport} &
Admissible link paths, inverse/composition laws, color representation,
rapidity scheme, and dynamical holonomy evaluation.\\
\texttt{GTMDOperator} &
Off-forward operator insertion and common Fock-overlap evaluator.\\
\texttt{ReductionMap} &
Typed TMD, GPD, current, form-factor, and OAM reductions with commuting
diagram tests.\\
\texttt{NuclearAmplitudeMap} &
Spin-1 Fock composition, coherent mechanisms, currents, and controlled
partial trace.\\
\texttt{PositiveCorrelator} &
Amplitude or \(LL^\dagger\) representation and irreducible TMD coordinates.\\
\texttt{TruncationTower} &
Regulator/basis/Fock filtration, matching maps, convergence fits, and
model-discrepancy estimates.\\
\texttt{ProvenanceComplex} &
Transition graph, no-double-counting audit, source/replacement information,
and observable ancestry.\\
\bottomrule
\end{longtable}

\subsection{Required property-based tests}

\begin{enumerate}
\item Intertwiners commute with every represented symmetry generator.
\item \(H_{\LF}=H_{\LF}^\dagger\), and its Poincare commutator defects decrease
along the truncation tower.
\item Wilson transport respects path concatenation and inversion.
\item Time reversal maps future and past link operators with the correct
T-even/T-odd behavior.
\item GTMD reduction diagrams commute within regulator and numerical
tolerances.
\item Density correlators remain positive without eigenvalue clipping.
\item Spin-1 rotations reproduce the irreducible \(0\oplus1\oplus2\)
transformation law.
\item Nuclear maps close total probability and momentum only after all
declared Fock sectors are included.
\item Disabling either amplitude in an OAM interference removes the
corresponding harmonic.
\item Sector-transition provenance detects deliberately duplicated pion,
gluon, or sea mechanisms.
\end{enumerate}

\subsection{Role of PennyLane}

PennyLane could prototype small symmetry-reduced Hamiltonian blocks or
state-composition circuits only after specifying the physical basis,
Hamiltonian, observables, and encoding.  A circuit should be accepted only
if it reproduces exact diagonalization in the same truncated Hilbert space
and offers a concrete advantage for variational eigensolving, differentiation,
or correlation diagnostics.  A generic entangling circuit with no map to
Fock amplitudes would add parameters without physics and should be rejected.

\section{Staged research program}

\subsection{Stage A: algebraic skeleton with analytic benchmarks}

\begin{enumerate}
\item Implement irreducible color, isospin, spin-1, and \(SO(2)\) OAM labels.
\item Encode the existing model's correlators as positive block tensors
without changing their central physics.
\item Implement path-groupoid link types and exact transformation laws.
\item Establish commuting reductions on analytic two- and three-body
light-front models.
\end{enumerate}
This stage tests the architecture without claiming new dynamics.

\subsection{Stage B: microscopic nucleon pilot}

Choose one controlled Hamiltonian route and solve at least
\[
|qqq\rangle\oplus|qqqg\rangle\oplus|qqqq\bar q\rangle
\]
at multiple truncations.  Compute a nonzero-\(\DT\) quark/gluon GTMD subset,
its TMD/GPD/current reductions, and OAM moments.  Compare with known nucleon
charges, PDFs, form factors, and lattice moments.  The pilot is successful
only if more than one observable family is predicted from shared
parameters.

\subsection{Stage C: dynamical gauge-link pilot}

Add an eikonal Wilson-line interaction derived from the same gauge sector.
Test:
\begin{itemize}
\item vanishing of naive-T-odd functions at zero rescattering;
\item exact future/past reversal;
\item flavor and operator differences from shared intermediate states;
\item separate gluon \(f\)- and \(d\)-type contractions;
\item stability against regulator and Fock enlargement.
\end{itemize}

\subsection{Stage D: spin-1 microscopic composition}

Construct the deuteron from a consistent \(NN\) interaction and current,
then add \(NN\pi\) and any retained non-nucleonic sectors with normalization
and subtraction rules.  Reproduce binding, static moments, elastic form
factors, and controlled \(b_1\) limits before predicting the full tensor TMD
sector.

\subsection{Stage E: inference, evolution, and withheld validation}

Match to a declared TMD scheme, evolve all ranks and coupled partonic
sectors, calibrate shared parameters, and reserve at least one observable
family for posterior predictive validation.  This stage is the first point
at which the phrase ``genuinely predictive next-level model'' could be
audited.

\section{Risks and explicit nonclaims}

\begin{longtable}{L{0.25\textwidth}L{0.34\textwidth}L{0.33\textwidth}}
\toprule
Risk & Failure mode & Required safeguard\\
\midrule
\endhead
Decorative topology &
Calling link classes or TMDs topological invariants without a dynamical
derivation. &
Restrict topology to path/transition organization; retain connection and
Hamiltonian dependence.\\
Symmetry overconstraint &
Exact isospin or truncated rotational symmetry erases physical breaking. &
Implement breaking operators and test the symmetric limit separately.\\
Tensor-network bias &
Low bond dimension suppresses long-range or OAM correlations. &
Publish bond-rank convergence alongside basis/Fock convergence.\\
Positivity misuse &
Demanding pointwise positivity of a Wigner quasidistribution. &
Apply convex positivity to the appropriate forward helicity correlator,
not to all phase-space representations.\\
Channel misuse &
Forcing coherent shadowing into an incoherent completely positive map. &
Retain amplitude phases and trace only after the coherent calculation.\\
Renormalization gap &
Treating a finite Hamiltonian matrix as QCD without counterterms or
matching. &
Specify regulator, counterterms, matching observables, and truncation flow.\\
Parameter migration &
Adding one latent coefficient per TMD inside sophisticated notation. &
Attach parameters only to microscopic vertices, currents, or controlled
discrepancy operators.\\
\bottomrule
\end{longtable}

\section{Recommended decision}

The recommended foundation is a
\emph{gauge-covariant graded light-front tensor architecture with
convex-state geometry}.  Its indispensable elements are:
\begin{enumerate}
\item the graded regulated light-front Hilbert space and Hamiltonian;
\item representation-theoretic intertwiners for color, flavor, spin, OAM,
and spin-1 polarization;
\item common GTMD operators and commuting observable reductions;
\item Wilson-line parallel transport represented on an admissible path
groupoid;
\item symplectic organization of transverse phase space and OAM;
\item positive-semidefinite amplitude geometry for forward correlators;
\item a filtered truncation tower with matching and convergence.
\end{enumerate}
The tensor-network implementation is strongly recommended as the sparse
computational realization.  The chain-complex layer is recommended initially
as a provenance and no-double-counting audit.  More ambitious topological
constructions should be deferred until a concrete physical observable or
computational obstruction requires them.

The proposal does not generate the interaction Hamiltonian, determine
counterterms, or solve QCD by terminology.  Its value is stricter and more
practical: it gives the microscopic calculation a structure in which flavor,
spin, OAM, gauge link, color, nuclear composition, positivity, and
uncertainty cannot drift apart.  If implemented faithfully, it would convert
many of the project's present validation tests into consequences of the
representation and would make the remaining failures scientifically
diagnostic.

\appendix
\section{Compact mathematical dictionary}

\begin{tabularx}{\textwidth}{L{0.29\textwidth}Y}
\toprule
Physical concept & Mathematical object\\
\midrule
Fock sectors & Grading and direct-sum Hilbert space\\
Momentum support & Simplex/constraint manifold \(\mathcal M_\nu\)\\
Flavor, color, spin, OAM & Group representations and irreducible labels\\
Allowed interaction & Equivariant intertwiner\\
Spin-1 vector/tensor polarization & Rank \(1\) and \(2\) spherical tensors\\
Sparse state contraction & Symmetry-preserving tensor network\\
Gauge link & Bundle parallel transport/holonomy\\
Process-dependent staple & Morphism in a path groupoid\\
GTMD & Gauge-covariant operator morphism/matrix element\\
TMD/GPD/current/OAM limits & Typed reduction maps\\
Consistency identity & Commuting diagram\\
Transverse OAM & Moment map on \(T^*\mathbb R^2\)\\
Forward positivity & Positive-semidefinite convex cone\\
Impulse nuclear dressing & Completely positive amplitude map where valid\\
Sector bookkeeping & Chain complex/provenance graph\\
Truncation sequence & Filtered Hilbert spaces/projective matching system\\
Predictive uncertainty & Pushforward of shared posterior plus discrepancy\\
\bottomrule
\end{tabularx}

\section{Acceptance questions for adopting the architecture}

\begin{enumerate}
\item Can every tensor index be mapped to a physical degree of freedom or
controlled numerical basis label?
\item Does every parameter belong to a Hamiltonian term, state amplitude,
current, matching coefficient, or explicit discrepancy operator?
\item Are Wilson paths and color representations part of operator identity?
\item Can at least two different observable reductions be computed from the
same nonzero-transfer correlator?
\item Is positivity structural rather than repaired?
\item Can exact isospin, pure-\(S\), zero-rescattering, and zero-additional-
Fock-sector limits be recovered?
\item Is every approximation embedded in a convergence sequence?
\item Does the architecture enable a genuinely withheld prediction?
\end{enumerate}
A negative answer does not necessarily reject the entire formulation, but
it identifies the next concrete design or physics problem.

\small
\begin{thebibliography}{99}

\bibitem{Vary2009}
J.~P.~Vary et al.,
``Hamiltonian light-front field theory in a basis function approach,''
Phys.\ Rev.\ C \textbf{81}, 035205 (2010),
arXiv:0905.1411.

\bibitem{Meissner2008}
S.~Mei{\ss}ner, A.~Metz, and M.~Schlegel,
``Generalized parton correlation functions for a spin-\(\tfrac12\)
hadron,'' JHEP \textbf{08}, 056 (2009),
arXiv:0906.5323; see also arXiv:0807.1154.

\bibitem{Lorce2012}
C.~Lorc\'e,
``Wilson lines and orbital angular momentum,''
Phys.\ Lett.\ B \textbf{719}, 185 (2013),
arXiv:1210.2581.

\bibitem{Cotogno2017}
S.~Cotogno, T.~van Daal, and P.~J.~Mulders,
``Positivity bounds on gluon TMDs for hadrons of spin \(\leq1\),''
JHEP \textbf{11}, 185 (2017),
arXiv:1709.07827.

\bibitem{BacchettaMulders2000}
A.~Bacchetta and P.~J.~Mulders,
``Deep inelastic leptoproduction of spin-one hadrons,''
Phys.\ Rev.\ D \textbf{62}, 114004 (2000),
arXiv:hep-ph/0007120.

\bibitem{Boer2016}
D.~Boer et al.,
``Gluon and Wilson loop TMDs for hadrons of spin \(\leq1\),''
JHEP \textbf{10}, 013 (2016),
arXiv:1607.01654.

\bibitem{Collins2011}
J.~Collins,
\emph{Foundations of Perturbative QCD},
Cambridge University Press (2011).

\end{thebibliography}

\end{document}
